\documentclass[a4paper,11pt]{article}
\synctex=1
\usepackage[margin=1in]{geometry}
%\documentclass[a4paper,sigconf,nonacm]{acmart}
%
\title{
Research Proposal:
Orbit-finitely Generated Algebraic Objects
}
\author{Arka Ghosh}%
\usepackage{amsmath,amssymb}
%\usepackage{mathpazo}
\usepackage[T1]{fontenc}
%
\usepackage[backend=biber,style=alphabetic,url=false,doi=false,isbn=false,eprint=false]{biblatex}
\setlength{\bibitemsep}{0pt}
\addbibresource{paper.bib}
\DeclareSourcemap{
  \maps[datatype=bibtex]{
    \map[overwrite]{
      \step[fieldset=editor, null]
      \step[fieldset=editortype, null]
      \step[fieldset=editora, null]
      \step[fieldset=editorb, null]
    }
  }
}
\AtEveryBibitem{
  \clearlist{location}
  \clearlist{address}
  \clearlist{publisher}
}
% !TEX root = research_plan.tex
%
\usepackage{hyperref}
\usepackage{xcolor}
%
\usepackage[justification=centering]{caption}
%
\usepackage{tikz}
\usetikzlibrary{arrows.meta,positioning,petri,cd}
%
\usepackage[framemethod=TikZ]{mdframed}
\usepackage{amsthm}%\usepackage[framed]{ntheorem}
\usepackage[capitalise]{cleveref}
\mdfdefinestyle{theorembox}{
    linecolor=black,
    linewidth=4pt,
    % --- This controls the "height" of the box around the text ---
    innertopmargin=0pt,     % Space at the top
    innerbottommargin=0pt,  % Space at the bottom
    innerleftmargin=0pt,   % Space on the left
    innerrightmargin=0pt,  % Space on the right
    % ------------------------------------------------------------
    roundcorner=0pt,        % Set to 5pt for rounded
    %backgroundcolor=gray!5, % Light background
    skipabove=0pt,%\baselineskip,
    skipbelow=0pt,%\baselineskip
}
\global\mdfdefinestyle{thm}{%
linecolor=black,
linewidth=1pt,
innertopmargin=0pt
}
\usepackage[framemethod=TikZ]{mdframed}
\usepackage{amsthm}
\usepackage{thmtools}
\usepackage[framemethod=TikZ]{mdframed}
\usepackage{amsthm}
\usepackage{thmtools}
\usepackage{xcolor}

% 1. Define the Frame appearance
\mdfdefinestyle{myThmFrame}{
    innertopmargin=4pt,
    innerbottommargin=4pt,
    innerleftmargin=10pt,
    innerrightmargin=10pt,
    linecolor=black,
    linewidth=1pt,
    %backgroundcolor=blue!5,
    roundcorner=0pt
}

% 2. Define the Theorem Style (The "All-in-one" definition)
\declaretheoremstyle[
    headfont=\bfseries,       % Bold title
    notefont=\itshape,        % Italicized optional note
    bodyfont=\normalfont,     % Standard text for content
    headpunct={. },           % Punctuation after "Theorem 1.1"
    postheadspace=0pt,        % MAGIC: 0pt forces text to stay on the same line
    mdframed={style=myThmFrame} % Attach the frame style here
]{framedrunin}

\declaretheorem[style=framedrunin, name=Theorem]{theorem}%[chapter]
\declaretheorem[style=framedrunin, name=Lemma,      sibling=theorem]{lemma}
\declaretheorem[style=framedrunin, name=Conjecture, sibling=theorem]{conjecture}
\declaretheorem[style=framedrunin, name=Question,   sibling=theorem]{question}
\declaretheorem[style=framedrunin, name=Note,       sibling=theorem]{note}
\declaretheorem[style=framedrunin, name=Example,    sibling=theorem]{example}
\declaretheorem[style=framedrunin, name=Definition, sibling=theorem]{definition}
\declaretheorem[style=framedrunin, name=Problem,    sibling=theorem]{problem}
\declaretheorem[style=framedrunin, name=Problem]{goal}

\newcommand\prob[4]{
  \begin{problem}[#1]\label{#2}
  $ $\\
  {\bfseries Input}: #3\\
  {\bfseries Question}: #4
  \end{problem}
}

%
%%% General text macros %%%
%
\newcounter{commentnumber}
\newcommand{\arka}[1]{\textbf{\textcolor{red}{Arka \stepcounter{commentnumber}\thecommentnumber: #1}}}
\newcommand{\mine}[1]{\textcolor{magenta}{\cite{#1}}}
\newcommand{\mythm}[2]{\textcolor{magenta}{\cite[#1]{#2}}}
\newcommand{\emphdef}[1]{\textcolor{purple}{\emph{#1}}}
\newcommand{\compsci}[1]{\textcolor{blue}{#1}}
\newcommand{\hlt}[1]{\textbf{\textcolor{violet}{#1}}}
%
%%% Domains %%%
%
\newcommand{\Q}{\mathbb{Q}}
\newcommand{\X}{\mathcal{X}}
\newcommand{\M}{\mathcal{M}}
\newcommand{\C}{\mathbb{C}}
\newcommand{\dloQ}{\mathcal{Q}}
\newcommand{\A}{\mathcal{A}}
\newcommand{\B}{\mathcal{B}}
\newcommand{\R}{\mathbb{R}}
\newcommand{\N}{\mathbb{N}}
\newcommand{\Z}{\mathbb{Z}}
\newcommand{\bitV}[2]{{\mathcal{V}_{#1}}{(#2)}}
\newcommand{\ff}{\mathbb{F}}
\newcommand{\fftwo}{\mathbb{F}_2}
\newcommand{\ordZ}{\mathcal{Z}}
\newcommand{\cU}{\mathcal{U}}
%
%%% General math macros
%
\newcommand{\setof}[2]{\left\{#1\ \mid\ #2 \right\}}
\newcommand{\pn}{\mathbf{P}}
%
%%% Atoms %%%
%
\newcommand{\aut}[1]{\mathsf{Aut}{\left(#1\right)}}
\newcommand{\age}[2][]{\textsc{Age}_{#1}{\left(#2\right)}}
\newcommand{\tup}[1]{\overline{#1}}
\newcommand{\otuequiv}[1]{\A^{(#1)}}
%\newcommand{\eq}[1]{\A^{(#1)}}
%
%%% Polynomials %%%
%
\newcommand{\poly}[2]{#1\left[#2\right]}
\newcommand{\idl}[1]{\mathcal{#1}}
\newcommand{\idlgen}[1]{\langle #1 \rangle}
%
%%% vector spaces %%%
%
\newcommand{\flin}[2][\Q]{#2 \to_{\mathsf{fin}} #1}
\newcommand{\bigd}[2][\Q]{#2 \to #1}
\newcommand{\defd}[2][\Q]{#2 \to_{\textsc{def}} #1}
\newcommand{\vgen}[1]{\langle #1 \rangle}
\newcommand{\ofspan}[2]{$\textsc{Span}_{#1}(#2)$}
\newcommand{\idvec}[1]{\mathbf{1}_{#1}}
%
%%% register automata %%%
%
\newcommand{\parikh}{\textsc{Parikh}}
%
%%% complexity %%%
%
\newcommand{\ptime}{\textsc{Ptime}}
\newcommand{\exptime}{\textsc{Exptime}}
\newcommand{\npc}{\textsc{NP-complete}}

\date{}
%%
\usepackage{titlesec}
\titlespacing*{\section}{0pt}{6pt}{3pt}
\titlespacing*{\subsection}{0pt}{6pt}{3pt}
\titlespacing*{\subsubsection}{0pt}{6pt}{3pt}
\titlespacing*{\paragraph}{0pt}{6pt}{3pt}
%
%
\begin{document}
%
%
%\let\allcaps=\rela
%
\maketitle
%
%\raggedright
\raggedbottom
%
\section{Introduction}
%
\paragraph*{Data-enriched models of computation.}
Modelling and verification are two fundamental aspects of computer science.
We model programs using abstractions such as finite automata or Petri nets,
%and then verify properties like non-emptiness or reachability automatically.
so we can verify their properties automatically.
For classical models over finite alphabets,
this approach has been highly successful,
supported by a mature theory and toolbox (for eg, SPIN).
In practice, however, many systems manipulate values drawn from infinite domains (timestamps, process identifiers, unbounded numbers, etc) so classical models are often too restrictive.
This motivates data-aware extensions such as register automata,
which extend finite automata with a finite set of registers for storing and comparing data values and accept data words (words where each position carries a label together with a data value),
and Petri nets with data, in which tokens carry data from an infinite data domain.
The data domains come in many forms.
Two common examples of data domains that appear in the literature are:
%
\begin{enumerate}
\item The equality domain, which is an infinite set of strings/identifiers/numbers that can be
checked against equality or disequality.
\item The ordered domain, which is the set of rational numbers with the usual ordering.
\end{enumerate}
These can be generalised to relational data domains,
where the elements of the data domain come from an infinite set with additional structures given by relations defined on it.
For equality and ordered data,
the relations are respectively equality and order.
Relational data domains can also be infinite graphs, where the elements of the domain are the nodes of the graph,
and the relation is the adjacency relation of the graph.
There are also algebraic data domains,
where instead of relations (like for relational data) one has functions/operations defined on them.
For example one can consider the domain of bit vectors \cite[Section 7.3.2]{atombook} which are arrays of binary bits with the bit-wise XOR operation.
One can also construct new data domains from old ones using simple operations such as disjoint sums and products.


While these data-enriched models better capture real systems, the added ability to store and compare data typically makes verification significantly harder.
A major source of this difficulty,
that we want to address in this project,
is the lack of algebraic foundations that,
in the infinite-data setting, play the same role that classical algebra plays for finite models.
In the finite case, algebraic methods are tightly linked to verification:
for example, Petri-net state equations can be formulated as integer linear programs,
and reachability for reversible Petri nets reduces to the polynomial ideal membership problem \cite{MM82}.
%
\paragraph*{Orbit-finite structures.}
The behaviour of data-enriched models is invariant under isomorphisms of the underlying data domain.
For example, a register automaton over the ordered domain that recognises the language of words that are sorted in the increasing order,
cares not about the specific letters but only about the ordering of consecutive letters.
Moreover, up to the isomorphism of the data domain, a data enriched model can behave only in finitely many ways.
For instance, a register automaton over the equality domain can only check if the incoming data value is equal or dis-equal to the data values stored in its registers.

In \cite{BKL14}, the intuitive understanding of finite up to isomorphism has been formalised by the notion of orbit-finiteness.
An \emphdef{orbit-finite set} is a set constructed from a data domain, that are finite up to symmetries (automorphisms) of the data domain.
The orbit-finite setting gives a uniform way to represent and analyse models of computation with infinite alphabets.
A register automaton is an automaton with orbit-finitely many states and transitions,
a data Petri net is a Petri net with orbit-finitely many places and transitions. This translation into the orbit-finite framework integrates naturally with algebra.
Transitions monoids of register automata are orbit-finitely generated,
state equations of data Petri nets are orbit-finite systems of linear equations,
the reachability problem for reversible data Petri nets reduces to the membership problem for orbit-finitely generated polynomial ideals.
%
\section{Research questions}
%
The overarching goal of this project is to \hlt{build a theory of orbit-finitely generated algebraic objects}.
In particular,
we aim to study three kinds of algebraic objects:
\emphdef{orbit-finite dimensional vector spaces},
\emphdef{polynomial rings with variables indexed by data values},
and \emphdef{orbit-finitely generated monoids}.
Given such an object, we are interested in its sub-objects (vector subspaces, polynomial ideals, sub-monoids) that are \emphdef{equivariant},
i.e.\ invariant under the symmetries of the underlying data domain.

We are interested in both structural properties and decision problems on these algebraic objects.
Structural properties usually come in one of the two following forms
%
\begin{enumerate}
\item \textbf{Finiteness properties:} Is every strictly increasing chain of equivariant vector subspaces/polynomial ideals finite? Does every orbit-finite monoid have finite J-height?
\item \textbf{Symmetry properties:} Does every equivariant system of polynomial equations have an equivariant solution? 
\end{enumerate}
%
The decision problems are usually one of the two following forms  
%
\begin{enumerate}
\item \textbf{Membership questions:} Decide whether a given vector (resp. polynomial) is inside the equivariant vector
subspace (resp. polynomial ideal) generated by another given finite set of vectors (resp. polynomials)?
\item \textbf{Finiteness questions:} Decide if a given orbit-finitely generated monoid is actually orbit-finite?
\end{enumerate}
%
Of course the answer to a specific question of the above kind depends on the underlying data domain.
Our aim is to \hlt{connect property of an algebraic object/decidability of an algebraic problem to structural properties of the underlying data domain}.
To do this we will study algebraic properties/problems for specific domains,
identify combinatorial properties responsible for satisfaction of an algebraic property/decidability of an algebraic problem, and extend the result from a specific domain to a general class.
This interplay between problems and properties leads to a connection with model theory.
%
\subsection{Connections with other research areas}
%
The origins of orbit-finite structures can be traced to the work of Fraenkel and Mostowski,
who introduced the \emph{set theory with atoms (ZFA)} and used \emph{permutation models} to show that the independence of axiom of choice in this setting
\cite{jech}.
It was rediscovered by Gabbay and Pitts as \emph{nominal sets},
which was used by them to provide a rigorous mathematical foundation for variable binding and $\alpha$-equivalence \cite{GP99}.
It was also independently rediscovered in the concurrency community in terms of \emph{named sets},
which was then used to give a compact representation for \emph{$\pi$-calculus processes}
\cite{MP98}.

Orbit-finitely generated algebraic objects have also been studied by several communities.
Polynomial rings with infinitely many variables were studied by Cohen in the context of metabelian groups \cite{Cohen67},
and by Aschenbrenner and Hillar for their applications in chemistry and algebraic statistics \cite{AH04}.
Orbit-finite dimensional vector spaces have been studied by David Evans because of their connection to Thomas conjecture \cite{CE91,homOFR25}.
%
%
%\section{Introduction}
%
%\paragraph*{Data-enriched models of computation.}
%Classical models of computation are extended with the capability of storing and manipulating data values from infinite domains to be able to model real world processes better.
%This leads to data-enriched models of computation such as register automata, which extend finite automata with a finite set of registers for storing and comparing data values and accept data words (words where each position carries a label together with a data value),
%and Petri nets with data, in which tokens carry data from an infinite data domain.
%
%The data domains come in many forms.
%The two most common data domains appearing in the literature are the \emphdef{equality domain} (an infinite set of strings/identifiers/numbers that can be checked against equality or dis-equality) and the \emphdef{ordered domain} (the set of rational numbers with the usual ordering).
%In general data domains can be more complex structures such as infinite graphs, hyper-graphs and trees.
%New data domains can be constructed from old ones by taking sums and products.
%
%%The class of data domains used in practice is quite diverse (\cite[Section 3.2]{cpnbook}).
%%In contrast, most of the literature in the theoretical computer science community focuses only on a limited set of data domains,
%%and usually on equality and ordered data.
%%\arka{verification difficult: we will tackle by studying relevant algebraic problems}
%
%%\arka{significance of the project: give a toolbox, but build algebraic foundation. orbit-finite math str, new combinatorial properties}
%
%%One of the reasons for this gap between theory and practice,
%%which I want to address in my project, is the lack of an algebraic foundation.
%While these data-enriched models better capture real world systems,
%the added ability to store and compare data typically makes their automatic analysis significantly harder.
%We aim to tackle this difficulty by studying algebraic problems useful in this context.
%In the finite case, algebraic methods are tightly linked to verification:
%for example, Petri net state equations are formulated as integer linear programs,
%reachability for reversible Petri nets reduces to the polynomial ideal membership problem \cite{MM82}.
%%
%%The behaviour of data-enriched models is invariant under isomorphisms of the underlying data domain.
%%Moreover, up to the isomorphism of the data domain, a data-enriched model can behave only in finitely many ways.
%%For example, a register automaton over the ordered domain that recognises the language of words whose letters that are sorted in the increasing order,
%%cares not about the specific letters but only about the ordering of consecutive letters,
%%and can only check if the incoming data value is smaller, bigger or equal to the data values stored in its registers.
%%
%%\vspace{-10pt}
%\paragraph*{Orbit-finiteness.}
%%
%The article \cite{BKL11} introduced the concept of \emphdef{orbit-finite sets}:
%sets constructed from a data domain, that are finite up to symmetries (symmetries) of the data domain.
%The orbit-finite setting gives a uniform way to represent and analyse data-enriched models of computation.
%A register automaton is an automaton with orbit-finitely many states and transitions,
%a data Petri net is a Petri net with orbit-finitely many places and transitions. This translation into the orbit-finite framework integrates naturally with algebra.
%State equations of data Petri nets are orbit-finite systems of linear equations, the reachability problem for reversible data Petri nets reduces to the membership problem for orbit-finitely generated polynomial ideals,
%transition monoids of register automata are orbit-finitely generated.
%%
%\section{Research questions}
%%
%The overarching goal of this project is to \hlt{build a theory of orbit-finitely generated algebraic objects}.
%In particular,
%we aim to study three kinds of algebraic objects:
%\emphdef{orbit-finite dimensional vector spaces},
%\emphdef{polynomial rings with variables indexed by data values},
%and \emphdef{orbit-finitely generated monoids}.
%Given such an object, we are interested in its sub-objects (vector subspaces, polynomial ideals, sub-monoids) that are \emphdef{equivariant},
%invariant under the symmetries of the underlying data domain.
%
%We are interested in both structural properties and decision problems on these algebraic objects.
%Structural properties usually come in one of the two following forms
%%
%\begin{enumerate}
%\item \textbf{Finiteness properties:} Is every strictly increasing chain of equivariant vector subspaces/polynomial ideals finite? Does every orbit-finite monoid have finite J-height?
%\item \textbf{Symmetry properties:} Does every equivariant system of polynomial equation have an equivariant solution? 
%\end{enumerate}
%%
%The decision problems are usually one of the two following forms  
%%
%\begin{enumerate}
%\item \textbf{Membership questions:} Decide whether a given vector (resp. polynomial) is inside the equivariant vector
%subspace (resp. polynomial ideal) generated by another given finite set of vectors (resp. polynomials)?
%\item \textbf{Finiteness questions:} Decide if a given orbit-finitely generated monoid is actually orbit-finite?
%\end{enumerate}
%%
%Of course the answer to a specific question of the above kind depends on the underlying data domain.
%Our aim is to \hlt{connect property of an algebraic object/decidability of an algebraic problem to structural properties of the underlying data domain}.
%To do this we will study algebraic properties/problems for specific domains,
%identify combinatorial properties responsible for satisfaction of an algebraic property/decidability of an algebraic problem, and extend the result from a specific domain to a general class.
%This interplay between problems and properties leads to a connection with model theory.
%%
%\tableofcontents
%
\section{State of the art and applicant's contribution}
%
%\arka{Is it better to have a separate state of the art or have one for each research question?}
%\arka{For LIP. have them together. And remove petri nets}
%
\paragraph*{Highlighting of citations.}
%
To highlight the contribution of the applicant to the state of the art,
citations to articles for which the applicant is a co-author (for eg, \mine{GL25}) are coloured in red.
%to separate them from the rest (for eg, \cite{atombook})
%\subsection{Research questions}
%
%\noindent
%The goal of this project is to lift classical algebraic results to the orbit-finite setting.
%In particular, we aim to study three kinds of algebraic objects:
%Three kinds of algebraic objects have been studied in the literature:
%\emphdef{orbit-finite dimensional vector spaces},
%\emphdef{polynomial rings with variables indexed by data values},
%and \emphdef{orbit-finitely generated monoids}.
%A sub-object (vector subspaces, polynomial ideals, sub-monoids) of such an algebraic object is called \emphdef{equivariant} if it is invariant under the symmetries of the underlying data domain.
%The questions on these objects can be divided into two broad categories.
%%
%\paragraph*{Finiteness properties.}
%Is every strictly increasing chain of equivariant vector subspaces (resp.\ polynomial ideals) finite?
%Does every orbit-finite monoid have finite J-height?
%Does every equivariant set of polynomials have a common root that is also equivariant?
%%
%\paragraph*{Membership problems.}
%%
%Can we decide whether a given vector (resp.\ polynomial) inside the equivariant vector subspace/cone (resp.\ polynomial ideal) generated by another given finite set of vectors (resp.\ polynomials)?
%Can we decide if a given orbit-finitely generated monoid is actually orbit-finite?
%\smallskip
%
%\arka{what about monoids?}
%
%\arka{structural prop, decision}
%We give a quick overview of the state of research on these objects.
%
\paragraph*{Orbit-finite dimensional vector spaces.}
These vector spaces were introduced in \cite{BKM21},
which showed that for equality and ordered data, strictly increasing chains of equivariant vector subspaces are finite,
and even gave an upper bound on their lengths.
This result was then used to prove decidability of the equivalence problem for weighted register automata.

Specific versions of orbit-finite systems of linear equations and inequalities constructed with the equality and ordered domain were studied by \cite{HLT17,HL18,HR22}.
In collaboration with his doctoral advisors, the applicant gave a complete picture of the solvability problem of orbit-finite systems of linear equations and inequalities constructed with the equality domain
\mine{GHL22,GHL25,arkathesis}.
%
%
%The articles \cite{HR22}, \mine{GHL22} and \mine{GHL25} studies orbit-finite systems of linear equations and inequalities, and
As a special case, this work shows that for the equality data domain,
the membership problems for finitely generated equivariant vector subspaces and cones are decidable.
As a corollary we get decidability of solvability of state equations of data Petri nets each of whose tokens carry a tuple of bounded length of elements from the equality domain.
Reachability problem for such data Petri nets is undecidable (\cite{Las16}).
Therefore solvability of state equations gives us a decidable over-approximation for reachability.
%
\paragraph*{Polynomial rings with variables indexed by data values.}
%
%
%\arka{inf mny inv under all permutations, gen by us, indexed by data values}
%
The polynomial ring with infinitely many variables has been studied in \cite{Cohen67},
which showed that this ring is Noetherian in the sense that every strictly increasing chain of ideals that are invariant under permutations of the variables, is finite.

%The polynomial ring with variables indexed by elements from the equality data domain was studied in \cite{Cohen67},
%which showed that this ring is Noetherian in the sense that every strictly increasing chain of equivariant ideals is finite.
This result was generalised in \mine{GL24} for polynomial rings with variables indexed by elements from an arbitrary data domain.
This paper showed that a polynomial ring such as this is Noetherian
(every strictly increasing chain of ideals that are invariant under symmetries of the underlying data domain, is finite)
exactly when the finite induced substructures of the data domain labelled with natural numbers are well-quasi-ordered under embeddability.
This condition is satisfied by both the equality and the ordered data domain.
As a corollary we get that for data domains satisfying the well-quasi-ordering condition,
strictly increasing chain of equivariant vector subspaces are finite and thus the equivalence problem for weighted automata with these data domain is decidable.
This generalises the result of \cite{BKM21}.

The article \mine{GL25} extended the results of \mine{GL24} by showing existence and computability of finite Gr\"{o}bner bases for the class of data domains that satisfy the well-quasi-ordering condition.
As a corollary we get decidability of the membership problem for orbit-finitely generated ideals for these data domains.
This in turn implies decidability of the reachability problem of reversible data Petri nets with the same data domains.
Here we must mention that the reachability problem for general data Petri nets is open for equality data and undecidable for ordered data.
Both of these data domains satisfy the well-quasi-ordering condition.


The results of \mine{GL25} also implies decidability of the membership problem of finitely generated vector subspaces of orbit-finite dimensional vector spaces over the data domains satisfying the well-quasi-ordering condition.
This generalises one of two main results of \mine{GHL22}.
%
\paragraph*{Orbit-finitely generated monoids.}
These have been introduced by \cite{nomMon},
which showed that when the syntactic monoid of the language of a register automaton (over equality data) is orbit-finite,
the language recognised by the register automaton is definable in first order logic if and only if the syntactic monoid is aperiodic.
This paper also showed that every orbit-finite monoid over the equality domain has finite \emph{$J$-height}.
%
\section{Detailed description of research questions}
%
\subsection{Orbit-finite dimensional vector spaces}\label{sec:ofvsp}
%
We start with an example of a coloured Petri net.
Through this example we explain the fundamental concept of \emph{orbit-finiteness}\footnote{A more detailed description can be found in \cite[Chapter 3]{atombook}.},
and motivate the study of \emph{orbit-finite dimensional vector spaces}.

Let \emphdef{$\dloQ$} denote the ordered domain,
i.e., the set of rational numbers with the usual ordering.
We write it as $\dloQ$ and not $\Q$ to emphasise that we consider it as a linear order and not a field.


Consider the coloured Petri net $\pn$ in \Cref{fig:PN},
which has a single place (drawn as a circle) and two transitions (drawn as rectangles),
and each of its tokens are labelled by a pair of elements from $\dloQ$.
%For brevity, we write a pair in $\dloQ^2$ as $ab$ instead of $(a,b)$.
%
\begin{figure}
\begin{minipage}[b]{0.55\textwidth}
\begin{center}
%\vspace{30pt}
\begin{tikzpicture}[
node distance=60pt,
on grid,>={Stealth[round]},
bend angle=45,auto,
every place/.style= {minimum size=15pt,thick},
every transition/.style={minimum size=15pt,thick}]
  \node[place]        (p) {};
  %\node[place,label=above:$q$,above=of p,structured tokens={$\#$}]        (q) {};
  \node[transition,right=of p,label=right:$a < b < c$] {$Y$}
    edge[pre,bend right]  node[above] {$(b,a)$} (p)
    edge[post, bend left] node[below] {$(b,c)$} (p);
  \node[transition,left=of p,label=left:$a < b < c$]   {$X$}
    edge[pre,bend left]    node[above] {$(a,b),(b,c)$} (p)
    edge[post, bend right] node[below] {$(c,a)$} (p);
\end{tikzpicture}
%%\vspace{5pt}
\caption{A Petri net $\pn$ with ordered data}\label{fig:PN}
%%\vspace{5pt}
\end{center}
\end{minipage}
\hfill
\begin{minipage}[b]{0.40\textwidth}
\begin{center}
%\vspace{10pt}
\begin{tikzcd}
(a_1,a_2) + (a_2,a_3) + (a_1,a_3)
\arrow[d,"{x(a_1,a_2,a_3)}"] \\
(a_3,a_1) + (a_1,a_3)
\arrow[d,"{y(a_1,a_3,a_4)}"] \\
(a_3,a_4) + (a_1,a_3)
%\arrow[d,"{x(a_1,a_3,a_4)}"] \\
%(a_4,a_1)
\end{tikzcd}
%\vspace{10pt}
\caption{A run of the Petri net $\pn$}\label{fig:run}
\end{center}
\end{minipage}
\end{figure}
%
%\arka{monoids will give space}

The transition $X$ can take two tokens labelled respectively by $(a,b)$ and $(b,c)$,
for any $a < b < c$,
and return a token labelled with $(c,a)$.
The transition $Y$ can take one token labelled with $(b,a)$,
for any $a < b$,
and return a token labelled with $(b,c)$ as long as $b < c$.
%Observe that the output of the transition $X$ is determined by the input.
%In contrast, the output of the transition $Y$ is only constrained by the input. 

A \emph{configuration} of the Petri net $\pn$ can be written as an element in the \emph{vector space freely generated by $\dloQ^2$}.
For example,
the configuration that contains $2$ tokens labelled by $(a_1,a_2)$ and one token labelled by $(a_2,a_3)$, for some $a_1 < a_2 < a_3$,
is written as $2(a_1,a_2) + (a_2,a_3)$\label{page:formal sum}.
The space of such vectors is denoted as \emph{$\flin{\dloQ^2}$},
as these vectors can also be interpreted as the set of functions from $\dloQ^2$ to $\Q$ which are non-zero only on finitely many elements of the domain.

Using this semantics,
each transition can be identified with a family of vectors,
and applying a transition can be thought as adding a vector from this family.
For example, for the Petri net $\pn$ we can write
\begin{equation}\label{eq:tran vec}
X = \setof{x(a,b,c)}{a < b < c\in\dloQ} \quad\text{and}\quad Y = \setof{y(a,b,c)}{a < b < c\in\dloQ}\ ,
\end{equation}
where $x(a,b,c) = (c,a) - (a,b) - (b,c)$ and $y(a,b,c) = (b,c) -  (b,a)$.
\Cref{fig:run} describes a run of $\pn$. 
In the figure $a_1 < a_2 < a_3 < a_4 \in\dloQ$.

\paragraph*{Orbit-finiteness.}
%Now we explain the concept of orbit-finiteness.
Let \emphdef{$\aut{\dloQ}$} denote the group of order preserving bijections of the set $\dloQ$.
It acts on the set $\dloQ^2$ in a point-wise manner.
That is,
for $\pi\in\aut{\dloQ}$ and $(a,b)\in\dloQ^2$,
$\pi((a,b)) = (\pi(a),\pi(b))$.
One can define an equivalence relation $\thicksim$ on $\dloQ^2$ as:
$(a,b) \thicksim (a',b')$ if $\pi((a,b)) = (a',b')$ for some $\pi\in\aut{\dloQ}$.
The equivalence classes are called \emphdef{orbits}.
The set $\dloQ^2$ is partitioned into finitely many orbits (three in particular), namely
\[
\setof{(a,b)\in\dloQ^2}{a < b}
\text{, }\setof{(a,b)\in\dloQ^2}{a > b}
\text{, and }
\setof{(a,b)\in\dloQ^2}{a = b}\ .
\]
%
In other words,
the set $\dloQ^2$ is \emphdef{orbit-finite}.
As a matter of fact, sets of the form $\dloQ^n$ are orbit-finite for every $n\in\N$.

%Observe that, the orbit of an element $(a_1,\dots,a_n)\in\dloQ^n$ can also be written as
%\[
%\setof{\pi((a_1,\dots,a_n))}{\pi\in\aut{\dloQ}} \ .
%\]
%\arka{Remove if unnecessary.}


In the general setting,
we consider a \emphdef{countable structure}
\label{page:str}
$\M = (M,R_1,\dots,R_m,f_1,\dots,f_n)$,
i.e.\ $M$ is a countably infinite set and $R_i$ and $f_j$ are respectively relations and functions on $M$ of finite arity.
By abuse of notation,
we use $\M^n$ to refer to the set $M^n$.
For example,
we write ``$a\in\M^n$ '', or in words ``$a$ is an element of $\M^n$ '' to mean that $a$ is an element of $M^n$.

We have already mentioned the structure of \emphdef{ordered domain $\dloQ = (Q,<)$} where $Q$ is the set of rational numbers and $<$ is the usual ordering;
and the structure of \emphdef{equality domain $\A = (A,=)$} where $A$ is an infinite set and $=$ is the equality relation.
Another example is the structure \emphdef{$\ordZ = (Z,<)$} where $Z$ is the set of integers and $<$ is the usual ordering.
%An example of a structure with functions and constants is a vector space \emphdef{$\bitV{\infty}{\fftwo} = (V,+,\mathbf{0})$} of infinite dimension over the finite field $\fftwo$ with two elements.
%Here $+$ denotes the usual addition of vectors and $\mathbf{0}$ denotes the zero vector.

An \emphdef{automorphism} of $\M$ is a bijection of its underlying set which preserves and reflects the structure.
We use \emphdef{$\aut{\M}$} to denote the set of automorphisms of $\M$.
For example:
\begin{enumerate}
  \item \emphdef{$\aut{\A}$} is the set of all bijections of $A$.
  \item \emphdef{$\aut{\dloQ}$} is the set of all order preserving bijections of $Q$.
  \item \emphdef{$\aut{\ordZ}$} is the set of all order preserving bijections of $Z$. Observe that every such bijection is a translation $n \mapsto n + d$ for some integer $d$.
  %\item \emphdef{$\aut{\bitV{\infty}{\fftwo}}$} is the set of all vector space isomorphism of $\bitV{\infty}{\fftwo}$.
\end{enumerate}
%
%A subset $B\subseteq \M^n$ is called \emphdef{equivariant} if for every vector $b$ in $B$ and automorphism $\pi$ of $\M$,
%$\pi(b)$ is also in $B$.
%
We will be interested in structures $\M$ for which $\M^n$ is orbit-finite for every $n$.
These structures are called \emphdef{oligomorphic}.
The structures $\A$ and $\dloQ$ are oligomorphic.
The structure $\ordZ$ is not oligomorphic since $\ordZ^2$ is not orbit-finite.
This is because sets of the form $\setof{(n,m)}{n-m=d}$ for integers $d$ form disjoint orbits.
Other examples of oligomorphic structures include the Rado graph, bit vectors and universal tree (\cite[Section 7]{atombook}).

We use \emphdef{$\flin{\M^n}$} to denote the set of functions from $\M^n$ to $\Q$ which are non-zero only on finitely many elements of the domain.
This set forms a vector space under pointwise addition and scalar multiplication.
The set $\M^n$,
seen as a subset of vectors in $\flin{\M^n}$,
forms a basis/generating set of $\flin{\M^n}$.
Thus, we write vectors in $\flin{\M^n}$ as finite formal sums of elements of $\M^n$.
As we are interested only in oligomorphic structures,
sets of the form $\M^n$ are orbit-finite.
Thus, we call vector spaces of the form $\flin{\M^n}$ to be \emphdef{orbit-finite dimensional}.
%
%The research on this vector spaces revolves around three topics,
%namely, equivariant vector spaces and linear Noetherianity,
%orbit-finite spanning problems, and
%the dual space of definable vectors.
%
%Vector spaces of this form has been studied mainly in two contexts.
%Namely, in the context of (weighted) register automata \cite{BFKM24},
%which has given has interesting properties such as linear Noethrianity (\Cref{subsec:equiv}) and orbit-finite spanning property (\Cref{subsec:dual});
%and in the context of coloured Petri nets \cite{HLT17,HL18},
%from where we get the decision problem of membership in a orbit-finitely spanned vector subspace (\Cref{subsec:span}).
%
%In the following sections we give an overview of the state of the art regarding these topics,
%and mention some concrete research directions that we will take in this project.
%
%\newpage
\subsubsection{Equivariant vector subspaces}\label{subsec:equiv}
%
The action of the group $\aut{\M}$ extends to $\flin{\M^n}$ as well.
For illustration,
for $\pi\in\aut{\dloQ}$ and $a < b < c$
\[
\pi(2(a,b) - (c,c)) =
2(\pi(a),\pi(b)) - (\pi(c),\pi(c)) \ .
\]
A vector subspace $V\subseteq \flin{\M^n}$ is called \emphdef{equivariant} if for every vector $v$ in $V$ and automorphism $\pi$ of $\M$,
the vector $\pi(v)$ is also in $V$.
For example,
the subspace of $\flin{\dloQ^2}$ spanned by $X\cup Y$ (\cref{eq:tran vec}) is equivariant.
%This subspace is also \emphdef{orbit-finitely spanned},
%i.e.\ spanned by an orbit-finite set.
%\arka{be consistent}
%In general,
%orbit-finitely spanned subspaces of $V$ are equivariant.
%However, the converse is not known to be true.
%
%\vspace{-2pt}
\begin{definition}\label{def:lin noe}
A structure $\M$ is called \emphdef{linearly Noetherian}
\label{page:lin noe}
if every strictly increasing chain of equivariant subspaces of vector spaces of the form $\flin{\M^n}$ is finite.
\end{definition}
%\vspace{-2pt}
%
Linear Noetherianity of a structure $\M$ implies \compsci{decidability of equivalence of weighted register automata} over $\M$ \mythm{Remark 33}{GL24}.

This property was first studied in \cite{BKM21} which showed that the structures $\A$ and $\dloQ$ are linearly Noetherian \cite[Lemma IV.8]{BKM21}.
In \cite[Page 21]{BKM21} the authors stated the following conjecture which we will be interested in this project.
%
%\vspace{-2pt}
\begin{conjecture}\label{conj:lin noeth}
Every oligomorphic structure is linearly Noetherian.
\end{conjecture}
%\arka{Jingjie's masters thesis?}
%\vspace{-10pt}
%
%In \cite{moreLen},
%this conjecture has been verified for smoothly approximated structures
%\cite[Page 440]{smooth},
%and for generically ordered expansions of Fra\"{i}sse limits with free amalgamation, in languages with at most binary relations \cite[Page 1602]{homsurvey}.
%%
%%the Rado graph\cite{wikiRado} and Henson graphs \cite{wikiHenson}.
%%It also follows 
%
%Linear Noetherianity of \emph{$\N$-nicely orderable} structures 
%follows from the results of both \cite{LS23} and \mine{GL24}.
%\arka{cite previous work}
%
\subsubsection{Orbit-finite spanning problems}\label{subsec:span}
%
A necessary condition for a configuration $v\in\flin{\dloQ^2}$ of $\pn$ to be reachable from another configuration $u\in\flin{\dloQ^2}$,
%
%\arka{tokens consumed and produced, can be written as equations}
%is that $(v - u)$ in the vector subspace $\vgen{X\cup Y}$ spanned by $X\cup Y$.
%This is an instance of the \emph{membership problem for orbit-finitely spanned vector subspaces},
%one of the decision problem we will focus on in my project.
%
%\arka{alternate}
%
is that there exists $z_1,\dots,z_n\in X\cup Y$ and $k_1,\dots,k_n\in\N$ such that
$(v - u) = \sum_{i=1}^n k_i\cdot z_i$.
This is an instance of an orbit-finite spanning problem, which we describe below.

For $K\subseteq\Q$,
a vector $x\in\flin{\M^n}$ is called \emphdef{$K$-spanned} by a set of vectors $P\subseteq\flin{\M^n}$ if there exists a vectors $p_1,\dots,p_n\in V$ and numbers $k_1,\dots,k_n\in K$ such that
$x = \sum_{i=1}^n k_i\cdot p_i$.
The \emphdef{orbit-finite $K$-spanning problem over $\M$} (in short \emphdef{\ofspan{K}{\M}}) is a decision problem of the following form:
\prob{\ofspan{K}{\M}}{prob:ofspan}{A vector $x\in\flin{\M^n}$, and an orbit-finite set of vectors $P\subseteq\flin{\M^n}$.}{Is $x$ $K$-spanned by $P$?}
%(\arka{must be on a single page})
%
We will work on \ofspan{K}{\M} for $K = \Q,\Q_+,\Z$ and $\N$.
These decision problem come up naturally from the study of \compsci{coloured Petri nets},
particularly in the study of the reachability problem and its relaxations.
\compsci{State equations of data Petri nets} are instances of \ofspan{\N}{\M},
and \ofspan{\Q_+}{\M} is useful for deciding \compsci{continuous reachability}.
%
%Orbit-finite spanning problems have been studied extensively for the equality domain in \cite{HR22} \mine{GHL22} and \mine{GHL25} .
%Specific versions of this problem for the ordered domain have been studied in \cite{HL18} and \mythm{Sections 4.3 and 4.5}{arkathesis}.
%Some of these results have been generalised by \mine{GL24,GL25} to the class of structures satisfying the well-quasi-order property mentioned in \Cref{sec:sota}.
%\arka{maybe discuss this in work strategy}
%\mythm{Theorem 6.1}{GHL22} shows decidability of \ofspan{\Q}{\A} and \ofspan{\Z}{\A}.
%The decidability of \ofspan{\N}{\A} follows from \mythm{Theorem 6.1}{GHL22} and \cite[Lemmas 11.12 and 11.14]{HR22}.
%Finally, the decidability of \ofspan{\Q_+}{\A} follows from \cite[Theorems 17 an 18]{GHL25}.
%
%In \mythm{Section 4.5}{arkathesis} we show decidability of a specific version of \ofspan{K}{\dloQ} for $K = \Q,\Z$.
%My result was soon generalised in two papers.
%In particular \cite[Theorem 1.1 and 2.8]{EvansRamsey} shows decidability of \ofspan{K}{\X} for $K=\Q,\Z$ for all structures $\X$ that admit a $\omega$-categorical Ramsey expansion.
%And \mythm{Corollary 40}{GL25} show decidability of \ofspan{K}{\X} for $K=\Q,\Z$ for all nicely orderable structures $\X$.
%\arka{cite previous work}
%
%A specific version of \ofspan{\Q_+}{\dloQ} and \ofspan{\N}{\dloQ} has been studied in \cite{HL18},
%where the vectors are from ${\flin{(\dloQ\times F)}}$ with $F$ being an arbitrary finite set.
%For this version, the $Q_+$-spanning problem is decidable in PTime,
%and the $\N$-spanning problem is Ackermann complete.
%
%\subsubsection{Dual Spaces}\label{subsec:dual}
%%
%%\arka{Write that it can be skipped}
%The \emphdef{dual} of a vector space $V$ is the set of all linear functions from $V$ to the underlying field,
%which forms a vector space under point-wise addition and scalar multiplication.
%Since the vector space $\flin{\M^n}$ is freely generated by $\M^n$,
%its dual is essentially the set $\bigd{\M^n}$ of all functions from $\M^n$ to $\Q$.
%This set also forms a vector space under point-wise addition and scalar multiplication. 
%
%This space is too big to be useful.
%In particular, it is uncountable, which makes it impossible to enumerate or even represent its elements.
%Both of these problems vanish if one restricts to vectors that are \emph{definable}.
%Intuitively,
%a function $\bigd{\M^n}$ is \emphdef{definable} if it can be defined using a first order formula which can use the relations of $\M$ and can use finitely many elements from $\M$ as constants.
%%
%\begin{example}\label{eg:def fun}
%For example,
%for $d \in \dloQ$, one can define a function $f_d : \bigd{\M^n}$ as
%\[
%\begin{aligned}
%f(d,a) &=  -1           &&\text{for $a < d$} &&\\
%f(d,a) &= \phantom{-}1  &&\text{for $a > d$} &&\\
%f(a,b) &= \phantom{-}0  &&\text{otherwise.}  &&
%\end{aligned} 
%\]
%\end{example}
%%
%We denote the set of definable functions in $\bigd{\M^n}$ as \emphdef{$\defd{\M^n}$}.
%As a consequence of oligmorphicity of $\M$,
%$\defd{\M^n}$ forms a vector space under point-wise addition and scalar multiplication.
%Intuitively, the sum of two functions $f$ and $f'$ in $\defd{\M^n}$, defined respectively using finite subsets $S$ and $S'$ of $\M$,
%can be defined using $S\cup S'$.
%
%Elements of $\defd{\M^n}$ can be used as a proof that a vector $x$ is not spanned by an orbit-finite set of vectors $P$.
%%
%\begin{example}\label{eg:def}
%Consider the set of vectors
%\[
%W = \setof{w(a,b,c)}{a < b < c\in\dloQ} \ ,
%\]
%where $w(a,b,c) = (b,a) + (b,c)$ for $a < b < c\in\dloQ$.
%Pick $d < d' \in \dloQ$.
%Say we want to check if the vector $(d,d') - (d',d)$ is $\Q$-spanned by $W$.
%Recall the function $f_d$ in \Cref{eg:def fun}.
%Extending $f_d$ linearly to $\flin{\dloQ^2}$ we get
%\[
%f_d((d,d') - (d',d)) = f_d(d,d') - f_d(d',d) = 1 \ ,
%\]
%and,
%\[
%f_d(w(a,b,c)) = f_d(b,a) + f_d(b,c) \ .
%\]
%If $b \neq d$, then $f_d(b,a) = f_d(b,c) = 0$,
%which implies $f_d(w(a,b,c)) = 0$ as well.
%And if $b = d$, then $f_d(d,a) = 1$ and $f_d(d,c) = -1$,
%making $f_d(w(a,d,c)) = 0$.
%Thus $f_d(w(a,b,c)) = 0$ for all $a < b < c$.
%This means $f_d(u) = 0$ for all vectors $u$ that are $\Q$-spanned by $W$.
%As $f_d((d,d') - (d',d)) \neq 0$,
%we prove it can not be $\Q$-spanned by $W$.
%\end{example}
%%
%We say $\M$ has the \emphdef{definable separator property} if for every vector $x\in\flin{\M^n}$ and vector subspace $U$ of $\flin{\M^n}$ such that $x\notin U$,
%there exists $f \in \defd{\M^n}$ such that $f(x) = 1$ and $f(U) = \{0\}$.
%This property has been shown to hold for oligomorphic Ramsey structures by \cite{EvansRamsey} and forms one of the core ingredients of the decidability of the $\Q$-spanning problem for structures admitting oligomorphic Ramsey extension.
%
%
%
%This property implies decidability of the $\Q$-spanning problem,
%as one can check if a vector $x$ is in the equivariant vector subspace spanned by an orbit-finite set of vectors $W$ by doing the following two things in parallel:
%\begin{enumerate}
%\item enumerate vectors that are $\Q$-spanned by $W$ and check if $x$ is one of them,
%\item enumerate functions in $\bigd{\M^n}$ and check if for any function $f$, $f(x)\neq 0$ and $f(W) = \{0\}$.
%\footnote{It is a general fact that definable functions can be evaluated on orbit-finite set \arka{citation?}.}
%\end{enumerate}
%
%
%
%
%
%
%%For example,
%%for every vector $x\in\flin{\M^n}$ and vector subspace $U$ of $\flin{\M^n}$ such that $x\notin U$,
%%there is an element $f$ in the dual such that $f(x) = 1$ and $f(U) = \{0\}$.
%%Thus, one can try to give an algorithm for checking if a vector $x$ is in the equivariant vector subspace spanned by an orbit-finite set of vectors $W$ by doing the following two things in parallel:
%%\begin{enumerate}
%%\item enumerate vectors that are $\Q$-spanned by $W$ and check if $x$ is one of them,
%%\item enumerate functions in $\bigd{\M^n}$ and check if for any function $f$, $f(x)\neq 0$ and $f(W) = \{0\}$.
%%\end{enumerate}
%%%
%%The set of vectors that are $\Q$-spanned by $W$ is countable,
%%therefore spanning them is not a problem.
%%However, the set $\bigd{\M^n}$ is uncountable, thus finding such an $f$ is not possible even by enumeration.
%%Moreover, it is unclear how one can represent such functions and evaluate them orbit-finite set of vectors.
%%
%%Both of these problems vanish if one restricts to vectors that are \emph{definable}.
%%Intuitively,
%%a function $\bigd{\M^n}$ is \emphdef{definable} if it can be defined using a first order formula which can use the relations of $\M$ and can use finitely many elements from $\M$ as constants.
%%%
%%
%%
%%We denote the set of definable functions in $\bigd{\M^n}$ as $\defd{\M^n}$.
%%As a consequence of oligmorphicity of $\M$,
%%the set of definable functions in $\bigd{\M^n}$ forms a vector space under point-wise addition and scalar multiplication.
%%Intuitively, the sum of functions $f$ and $f'$, defined respectively using finite subsets $S$ and $S'$ of $\M$,
%%can be defined using $S\cup S'$.
%%
%%As illustrated in \Cref{eg:def}, a definable function can be evaluated on an orbit-finite set of vectors.
%%This is an instance of general fact that definable functions can be evaluated on orbit-finite sets.\arka{citation?}
%%
%%In my project we would be interested in the following two properties of $\defd{\M^n}$.
%%
%%We say $\M$ has the \emphdef{definable separator property} if  
%
%We will also be interested in the \emphdef{orbit-finite spanning(basis) property} of a structure $\M$:
%the space $\defd{\M^n}$ has an orbit-finite spanning set(basis) for every $n$.
%% 
%\begin{example}\label{def:basis}
%Consider the structure $\dloQ$.
%The only subsets of $\dloQ$ that can be defined using the order $<$ are finite unions of points of intervals.
%This implies every definable function $\defd{\dloQ}$ essentially is a partition of $\dloQ$ into finitely many points and intervals,
%each of which is coloured by a number from $\Q$.
%This means,
%$\defd{\dloQ}$ is spanned by the following four orbits of vectors
%%
%\begin{align*}
%&\setof{\idvec{\{a\}}\hspace{12pt}}{a\in\dloQ}\\
%&\setof{\idvec{(a,b)}\hspace{7.4pt}}{a < b\in\dloQ}\\
%&\setof{\idvec{(a,+\infty)\hspace{1pt}}}{a\in\dloQ}\\
%&\setof{\idvec{(-\infty,a)}}{a\in\dloQ}\ .\\
%\end{align*}
%%
%\end{example}
%%
%Theorem 6.7 in \cite{BFKM24} says that $\A$ and $\dloQ$ satisfies the orbit-finite spanning property.
%This result was useful for their proof of decidability of \compsci{equivalence of weighted/unambiguous register automata}.
%All of the articles that studied the equivalence problem before \cite{BFKM24}, studied it with the \emph{non-guessing}
%\footnote{
%A register automata is non-guessing if the letters stored in its register at every point must come from the part of the word that it has read at that point.}
%restriction.
%
%
%
%
%Part of the result of \cite{BFKM24} was improved by \mythm{Theorem 3.3}{GHL22} where we show that $\A$ has the orbit-finite basis property.
%These results were improved significantly by \cite{Prz24} who showed that all oligomorphic and $\omega$-categorical structures, and also the structure $\dloQ$,
%has the orbit-finite basis property \cite[Theorem 3.7 and 3.11]{Prz24}.
%The author also conjectures that the orbit-finite basis property is equivalent to being oligomorphic and NIP \cite[Conjecture 3.1]{Prz24}.
%
%As a long term objective,
%we would also like to study the orbit-finite basis property \mythm{Question 8.2}{arkathesis}.
%This has been proven for oligomorphic and $\omega$-stable structures \cite[Theorem 3.7]{Prz24} and the ordered domain \cite[Theorem 3.11]{Prz24},
%and has been conjectured to be equivalent to being NIP for oligomorphic structures \cite[Conjecture 3.1]{Prz24}.
%A short term objective towards settling this conjecture is to study this property for specific structures
%
%
%\subsection{Polynomial rings with infinitely many variables}
%
%\subsection{Orbit-finitely generated monoids}
%\arka{Example or not?}
%\arka{See how large it makes the draft}
%
%
%\subsection{Polynomial rings with infinitely many variables}
%
%\subsection{Orbit-finitely generated monoids}
%
%\newpage
%
\subsection{Polynomial rings with variables indexed by data values}\label{sec:poly}
%
In this setting, we study polynomial rings $\poly{\Q}{X_{\M}}$ with rational coefficients variables indexed by elements from a data domain $\M$.
Automorphisms of $\M$ act on polynomials in a variable-wise manner.
For example,
for an automorphism $\pi$ of $\M$ and $a,b,c\in\M$,
\[
\pi(x_a - 3x_b^2x_c + 4) = x_{\pi(a)} - 3x_{\pi(b)}^2x_{\pi(c)} + 4
\]
An \emphdef{equivariant ideal} in such a polynomial ring is an ideal that is invariant under the renaming of variables.
For instance, the set of polynomials whose coefficients sum up to zero is an equivariant ideal in the polynomial ring with elements from the equality domain as variables.

The three cornerstones of the theory of polynomial rings with finitely many variables are:
\begin{enumerate}
\item \textbf{Hilbert's basis theorem}: Every ideal has a finite set of generators.
\item \textbf{Gr\"{o}bner bases:} These are specific kind of generators for an ideal with good computational properties.
Given a set of generators $G$ for an ideal,
one can use the \textbf{Buchbergers algorithm} to compute a Gr\"{o}bner basis $B$ of the ideal generated by $G$.
Then, one can check whether another polynomial $f$ is in the ideal generated by $G$ (equivalently, $B$) by reducing $f$ by $B$ using the multivariate division algorithm.
\item \textbf{Hilbert's nullstellensatz:} A finite set of polynomials has a common complex root if and only if the ideal generated by the polynomials does not contain the unit polynomial $1$.
\end{enumerate}
%
In \mine{GL24} we showed that every equivariant ideal in $\poly{\Q}{X_{\M}}$ has a finite set of generators if and only if finite substructures of $\M$ labelled with natural numbers, are well-quasi-ordered under embeddability.
In \mine{GL25} we extended this result by showing that for the same class of structures,
every equivariant ideal has an orbit-finite Gr\"{o}bener basis which can be computed using the equivariant extension of the Buchberger's algorithm.

Serge Lang has shown in \cite{Lang52} that every countable system of polynomial equations with infinitely many variables has a complex solution if and only if the ideal generated by the polynomials does not contain the unit polynomial $1$.
In this project we will study a variation of this result in the orbit-finite setting.
\begin{question}
Characterise data domains $\M$ satisfying the equivariant nullstellensatz property:
every orbit-finite set of polynomials $G$ in $\poly{\Q}{X_{\M}}$ that has a common root,
has a common equivariant root.
\end{question}
%
Interestingly this property is true for a large class of structures when the underlying field is finite \cite{Prz20}.%
\subsection{Orbit-finitely generated monoids}\label{sec:ofmon}
%
%\arka{add example?}
The article \cite{nomMon} conjectured that the Sch\"{u}tzenberger-McNaughton-Papert theorem extends to register automata:
\begin{conjecture}
The language recognised by a register automaton is definable in first order logic if and only if its syntactic monoid is aperiodic.
\end{conjecture}
The same article proved it for the case when the syntactic monoid is orbit-finite.
In this project we will be interested in the settling the full version of the conjecture.

The syntactic monoid of the language of a register automata is orbit-finitely generated,
but not always orbit-finite.
In fact, orbit-finite monoids recognise the same class of languages as single-use register automata \cite[Theorem 6]{BS20},
which is a strictly weaker model.

As single-use register automata are better behaved than the general class of register automata,
we would be interested in the following problem:
%of deciding whether the syntactic monoid of a given register automaton is orbit-finite.
%
\prob{Finiteness of syntactic monoid}{prob:ofmon}{A register automaton.}{Is the syntactic monoid of the language of the register automaton orbit-finite?}
%
This problem has shown to be decidable for register automata over equality domain by \cite{ofMon}.
However, their methods do not extend to other data domains such as the ordered domain.

The decidability of this problem is related to the finiteness of the $J$-height  \cite[Page 5]{Colcombet11}.
%\arka{say a bit more about $J$-height}
Lemma 9.3 in \cite{nomMon} says that every orbit-finite monoid over the equality domain has finite \emph{$J$-height}.
However the same does not hold for the ordered domain.
This gives two questions that we would be interested in:
\begin{question}
For which data domains orbit-finite monoids have finite $J$-height?
Can the definition of $J$-height be modified so that orbit-finite monoids over any oligomorphic data domain have finite $J$-height?
\end{question}
%\newpage
\section{Organisational details of the project}
%
\subsection{Research methodology}
%
%\arka{shorten. Maybe not too much since also used in risk assesment}
Our ultimate objective is to characterise properties of algebraic objects and decidability of algebraic problems by structural properties of the underlying data domain.
An example of such a characterisation are Theorems 11 and 12 in \mine{GL24},
which characterise the Noetherian property of polynomial rings with elements from a data domain as variables by a well-quasi-ordering property of the data domain.

Our strategy to achieve this objective is to study an algebraic property/problem for a specific data domain,
find a structural property of the data domain that leads to satisfaction of the algebraic property/(un)decidability of the algebraic problem,
find another data domain not having this structural property,
and continue.
%
%What we call data domains, is called \emph{structures} by model theorists. 
%Classification of structures is the objective of model theory.
%Towards these objective, model theorists have identified several properties of structures.
%These properties come in forms.
%For instance,
%in terms of amalgamation property of their finite induced substructures (free amalgamation, strong amalgamation),
%in terms of the action of their automorphism groups (extreme amenability),
%and finally,
%from the classification theory of Shelah (stability and independence).
%We will try to connect this well-known properties with properties of algebraic objects/decidability of algebraic problems.
%
\subsection{Dissemination of the results}
%
In theoretical computer science it is more common to publish in peer reviewed conference proceedings in comparison to journals.
Following this custom,
we will submit our results in conferences on logic and verification (LICS, ICALP, CAV, FOSSACS, CONCUR, MFCS,...).
Full versions of the conference papers will be submitted to relevant journals (Journal of the ACM, LMCS, Theoretical computer science,...).
Preprints of the publications will be made freely available through arxiv.
The results will be popularised by presenting them in relevant conferences and research seminars.

Finally, we will build a python library where we will implement our algorithms.
This will help us to showcase our results and at the same time benchmark our algorithms.
%
\subsection{Project timeline}
%
The research on the three kinds of algebraic objects (orbit-finite dimensional vector spaces, polynomial rings with variables indexed by data values, and orbit-finitely generated monoids) can be carried out independently.
For all these objects,
we will first study them for specific data domains,
followed by their study on general classes of data domains.

We will start our investigation with polynomial rings,
as this will be continuation of our ongoing research.
This will be followed by vector spaces,
since the applicant has gained prior experience in this topic during his doctoral studies.
Monoids will be studied towards the end of the project.
%
\subsection{Description of the team}
%
The applicant is the sole member of the research team.
%
\subsection{Timeliness of the project}
%
In the past few decades, the study of data enriched models of computation has caught the attention of the theoretical computer science community
(\cite{KF94},\cite{LNORW08},\cite{atombook}).
With it, has started the study of orbit-finitely generated algebraic objects
(\cite{BKM21},\mine{GL24},\cite{nomMon}).
However, most of the literature till now only studied these algebraic objects only for specific data domains such as equality and ordered.
This makes our project,
whose aim is to study them for general classes of data domains,
extremely timely.
%
\subsection{Risk analysis}
%
Our aim is to connect property of an algebraic object/decidability of an algebraic problem to structural properties of the underlying data domain.
This is an ambitious goal.
However, we have been successful in achieving this objective for existence and computability of Gr\"{o}bner bases for equivariant polynomial ideals (\mine{GL24,GL25}).
%\footnote{\arka{variables indexed by data values? orbit-finitely many variables}}
So it is realistic to expect further results in the same spirit.

Our ambitious goal is combined with achievable milestones.
Even if we do not see much success in the investigation of algebraic objects over general classes of data domains,
we expect significant success in the study of these objects over specific data domains (equality data, ordered data, Rado graph,...). 
%
\subsection{Local collaboration}
%
We will collaborate with Samuel Braunfeld,
who is a researcher at the Computer Science Institute of the Czech Academy of Science, working in model theory and combinatorics.
Due to the connection of this project with model theory,
his expertise will be very useful.
%
\subsection{International collaboration}
%
Orbit-finite algebraic objects is an active area of research among several international scholars.
Currently, the applicant is collaborating with Aliaume Lopez and Joanna Fijalkow,
both of whom are members of his current institution (\textbf{LaBRI, Bordeaux}).
These collaborations will be continued in the project. 

We expect collaboration with Bartek Klin at the \textbf{University of Oxford}, and his PhD student Jingjie Yang.
Both of them have made fundamental contributions to the study of orbit-finite dimensional vector spaces.

Finally, we will also collaborate with Miko\l{}aj Boja\'{n}czyk, S\l{}awomir Lasota and Piotr Hofman at the \textbf{University of Warsaw}.
Alongside Bartek Klin, Miko\l{}aj Boja\'{n}czyk and S\l{}awomir Lasota are the pioneers of the study of orbit-finite sets and structures.
Piotr Hofman has made significant contributions in the study of orbit-finite spanning problems.
%
\subsection{Travel plans}
%
We expect one or two publications in peer-reviewed conference proceedings every year.
Therefore, we expect two conference travels per year on average with active participation.

Apart from conference travels, the applicant will also make two research visits per year.
The first one or two visits will be to LaBRI.
Subsequent visits will be made in accordance with established collaboration.
In addition to collaboration on the core problems of the project,
the applicant also present the results of the projects in research seminars during these visits.

%
\subsection{Potential benefits and future applications}
%
This project will benefit the study of verification problems for data enriched models of computation in both direct and indirect manners.


The class of data domains used in practice is quite diverse (\cite[Section 3.2]{cpnbook}).
By contrast, much of the theoretical computer science literature concentrates on a relatively small family of data domains,
and usually on equality and ordered data.
Our project will develop an algebraic toolbox to support the study of verification problems for data-enriched models of computation over general classes of data domains.

At the same time, we expect to relate well-known model-theoretic properties of data domains to algorithmic properties.
Establishing these connections will strengthen verification methods for data-enriched models of computation.
%
\subsection{Data management}
%
The project does not involve or require collection of any kind of data.
%
\subsection{Gender equity plan}
%
We will adhere to the principles of the Gender Equality Plan (\url{https://www.cs.cas.cz/docs/ENG_Plan_gender.pdf}) of the Institute of Computer Science of Czech Academy of Sciences.
%
\printbibliography
%
\end{document}
